\section{Introduction}
\lipsum[1]
This section also demonstrates a citation to prior work~\cite{lamport1994latex,knuth1984texbook,shannon1948mathematical,he2016deep} and a reference to \sect{example}.

\section{Example section}\llabel{sec:example}
% llabel{} is a chapter-aware \label{} that does not conflict with labels in other chapters.
% Use \sect{}, \fig{}, \tbl{}, and \algo{} to refer to such labels from within the same chapter.
\lipsum[2]

\subsection{Example equation}
\begin{equation}
  y(t) = \sum_{k=1}^{N} a_k\, x(t-k) + \epsilon(t),
  \label{eq:chapter-a-example}
\end{equation}
where $a_k$ are coefficients, $x(t)$ is the input, and $\epsilon(t)$ is noise, as in Eq.~\eqref{eq:chapter-a-example}.

\subsection{Example figure}
\fig{example} shows a placeholder figure. Replace the framed box with \verb|\includegraphics| and a file path.

\begin{figure}[htbp]
  \centering
  \framebox[0.6\textwidth]{\parbox[c][4cm][c]{0.55\textwidth}{\centering Placeholder for figure}}
  % \includegraphics[width=0.6\textwidth]{figures/example.pdf}
  \caption{Example figure caption.}
  \llabel{fig:example}
\end{figure}

\subsection{Example table}
\tbl{example} shows a placeholder table.

\begin{table}[htbp]
  \centering
  \caption{Example table caption.}
  \llabel{tbl:example}
  \begin{tabular}{lcc}
    \hline
    Method   & Metric 1 & Metric 2 \\
    \hline
    Baseline & 0.80     & 0.70     \\
    Proposed & 0.90     & 0.85     \\
    \hline
  \end{tabular}
\end{table}

\subsection{Example algorithm}
\algo{example} shows a placeholder algorithm.

\begin{algorithm}[htbp]
  \caption{Example algorithm}
  \llabel{algo:example}
  \begin{algorithmic}[1]
    \Require Input $x$
    \Ensure Output $y$
    \State $y \gets 0$
    \For{$i = 1$ to $N$}
      \State $y \gets y + x_i$
    \EndFor
    \State \Return $y$
  \end{algorithmic}
\end{algorithm}

\section{Summary}
\lipsum[3]
